\section{题目模板}

\begin{question}{重要公式}
	123
\end{question}

\section{线性方程组}

\subsection{方程组概念练习}

\begin{enumerate}
	% --- 判断题 ---
	\item 若两个方程组是同解的，则这两个方程组中方程个数必一样多。
	\begin{choices}
		\item 正确
		\item \textbf{错误}
	\end{choices}
	
	\item 若两个方程组有相同的解，就称这两个方程组是同解的。
	\begin{choices}
		\item 正确
		\item \textbf{错误}
	\end{choices}
	
	\item 若 $n$ 维有序数组使方程组中的每个方程成为恒等式，则称这个数组就是方程组的一个解。
	\begin{choices}
		\item \textbf{正确}
		\item 错误
	\end{choices}
	
	\item 若 $n$ 维有序数组使方程组中的某个方程成为恒等式，则称这个数组就是方程组的一个解。
	\begin{choices}
		\item 正确
		\item \textbf{错误}
	\end{choices}
	
	% --- 多选题 ---
	\item 对于线性方程组，主要研究方程组的哪些问题？（多选）
	\begin{choices}
		\item \textbf{线性方程组是否有解？}
		\item \textbf{若线性方程组有解，其解有多少个？如何求解？}
		\item \textbf{若线性方程组有解，其解的一般形式如何？或其解集的结构如何？}
		\item 线性方程组如何用数值方法求解？
	\end{choices}
	
	\item 以下属于线性方程组的是：
	\begin{choices}
		\item $\begin{cases} a_{11}x_1 + a_{12}x_2 + \dots + a_{1n}x_n = b_1 \\ \dots \\ a_{m1}x_1 + a_{m2}x_2 + \dots + a_{mn}x_n = b_m \end{cases}$
		\item 包含变量乘积项的方程组（如 $x_1x_2=1$）
		\item 包含变量指数项的方程组（如 $x_1^2=1$）
		\item 包含三角函数的方程组（如 $\sin x_1=0$）
	\end{choices}
\end{enumerate}